% ICLR 2027 end-of-main-text statements. All three sit AFTER the discussion and
% BEFORE the references, and NONE of them count toward the 9-page limit.
%   - AI use statement      : REQUIRED by the ICLR 2027 AI Policy for Authors, max 1 page
%   - Ethics statement      : recommended, max 1 page
%   - Reproducibility stmt  : recommended
% Use \subsection*/\section* (unnumbered) so they do not disturb section numbering.

\section*{AI use statement}
% ============================ DRAFT — AUTHOR REVIEW REQUIRED ============================
% This is a first draft written from the actual record of how this project was run.
% Verify every sentence against your own recollection before submitting: you are the
% one taking responsibility for it. Trim anything that does not apply.
% =======================================================================================
\TODO{Author: review and edit this statement before submission.}
We used a generative AI coding assistant extensively in the \emph{execution} of this
work: writing and debugging experiment, evaluation and data-preparation code,
launching and monitoring cluster jobs, aggregating metrics from raw per-sample
result dumps, and producing diagnostic figures. We also used it as a discussion
partner when interpreting results and when deciding which experiments to run
next. All research questions, experimental designs, and claims in this paper were
set by the authors, and every quantitative result reported here was recomputed by
the authors from the on-disk artifacts produced by the evaluation runs.

We did not use generative AI to fabricate, extrapolate, or impute any experimental
number, nor to generate any figure content that is not a direct rendering of
measured data. All prose in this paper was written or rewritten by the authors.

We highlight two cases where AI-assisted analysis produced an error that the
authors' verification caught, because they bear on how the results should be
read. First, an offline re-analysis of our anomaly maps initially reported a gain
from replacing the max-pooled image score with a high percentile; this turned out
to be an artifact of nearest-neighbour upsampling in the analysis script, and the
effect reverses under the bilinear interpolation the model actually uses. Second,
an early causal explanation of our MPDD failure mode was retracted after the
intervention it implied was measured and made results worse. Both are reported as
negative results in this paper. We take responsibility for the final content of
this work, including all text, claims, and artifacts.

\section*{Ethics statement}
\TODO{Author: confirm or delete.}
This work studies defect detection on public industrial anomaly-detection
benchmarks (MVTec-AD, VisA, MPDD, Real-IAD). It involves no human subjects, no
personally identifying information, and no new data collection. The datasets are
used under their published licences and for their intended research purpose. We
note one deployment consideration: the image-level detector reported here is not
calibrated for a specific operating point, and our per-object results show that
accuracy varies substantially across object classes, so it should not be treated
as drop-in quality control without per-class validation.

\section*{Reproducibility statement}
\TODO{Author: confirm the code/artifact release plan, then delete this note.}
All models are trained on a single public dataset (MVTec-AD) and evaluated
zero-shot on the others; no test-set data enters training or calibration. The
training recipe, LoRA configuration, loss and readout hyper-parameters, and the
frozen image-score calibration constants are given in \cref{app:impl}, and the
construction of the instruction data in \cref{app:data}. The evaluation protocol,
including how images with no segmentation readout are handled, is specified in
\cref{app:eval}. Reported aggregates are means over object categories; the seed
sensitivity of each metric is quantified in \cref{app:ablation}, and all
configuration comparisons are reported against those bands. We release the
training and evaluation code together with the per-sample result dumps from which
every table in this paper can be recomputed.
